\documentclass{article}
\usepackage[T1]{fontenc}
\usepackage{lmodern}
\usepackage{graphicx}
\usepackage{amsmath,amssymb}
\usepackage[a4paper,margin=2cm]{geometry}
\usepackage[absolute,overlay]{textpos}
\setlength{\TPHorizModule}{1mm}
\setlength{\TPVertModule}{1mm}
\usepackage[indonesian]{babel}
\usepackage{hyperref}
\setlength{\emergencystretch}{2em}
% unit: Halaman 15

\begin{document}

% === Halaman 15 (python/native) ===

15

Alice melakukan enkripsi dengan kunci public Bob (y = 1185, g = 2, p = 2357):


a = gk mod p = 21520 mod 2357 = 1430


b = ykm mod p = 11851520  2035 mod 2357 = 697

 Jadi, cipherteks yang dihasilkan adalah (1430, 697). 

 Alice mengirim cipherteks ini kepada Bob.

(c) Dekripsi (Oleh Bob)

  Bob melakukan dekripsi dengan kunci privatnya (x = 1751, p = 2357):


(ax)–1 = ap – 1 – x mod p  = 14302357 – 1 – 1751 mod 2357 = 1430605 mod 2357 = 872


m = b/ax mod p  = b  (ax)–1 mod p = 697  872 mod 2357 = 2035

     Bob mendapatkan kembali plainteks m = 2035 yang dikirim oleh Alice.

\end{document}
